% MORSE motivation draft (2026-10-15). Drop-in text for morse_cvpr_revised.tex, written in the strong version:
% every [X] is filled from the analyses listed at the end of this file (run them with the revised joint groups,
% v8_joint_groups_rev_k8.json). Notation and labels follow the draft; no new bib entries are needed.
%
% Part A replaces Introduction paragraphs 2-4 and the contribution list (paragraph 1 and the teaser figure stay).
% Part B is a new section between Problem Formulation (Sec. 3) and Method (Sec. 4).

% =====================================================================================================================
% PART A: INTRODUCTION (after the unchanged first paragraph)
% =====================================================================================================================

A widely used design for such policies handles robots with different numbers of joints by processing every joint
with shared modules and predicting the action of every joint with a shared decoder, conditioned on a description of
the joint and the body~\citep{huang2020smp,gupta2022metamorph,bohlinger2024onepolicy}. This design lets one network
control any body, but it also fixes how the robots share it: every joint of every robot updates the same parameters.
The hip-abduction joint of a quadruped, the knee of a humanoid, and the coxa of a hexapod rotate about different axes,
carry different loads, and play different roles in the gait, and yet their updates accumulate on the same weights;
conditioning on the descriptor changes the input of the shared modules, not the parameters they update. We find that
this sharing makes the updates of dissimilar joints interfere: in a trained policy, the gradients of joints of
different types conflict in [X]\% of pairs, even within one robot, whereas joints of the same type agree across
robots (Sec.~\ref{sec:where}).

Existing remedies change how the robots update this shared network rather than what they share. Gradient surgery
projects conflicting gradients onto each other~\citep{yu2020pcgrad}, and task grouping updates groups of tasks in
turn~\citep{jeong2025sel}. Embodiment Grouping (EG) brings the latter to robots: it clusters them by whole-body
morphology and updates the policy one group at a time, motivated by the conflicts between the gradients of different
robots~\citep{abe2026crossembodimentoffline}. A grouped update, however, changes when a robot updates the parameters,
not which parameters it updates. Every group still writes to every weight of the policy, and the optimizer state that
the groups share carries the step of one group into the next. Sequential group updates have been shown to help by
improving the learning of task-specific parameters~\citep{jeong2025sel}, which a fully shared policy does not have,
and indeed, conflicting gradients are as frequent within an EG group as between groups [X] (Sec.~\ref{sec:where}).
Moreover, the unit of all these methods is the whole robot, whereas robots are alike in some joints and different in
others: a decision made per robot can neither separate the dissimilar joints of similar robots nor join the similar
joints of dissimilar ones.

These observations point to two requirements for sharing in a cross-embodiment policy. First, interfering updates
should be assigned to different parameters, rather than to the same parameters at different times. Second, the unit
that decides the sharing should match the unit of the computation. The joint encoder and the decoder process one joint
at a time, and their gradients agree between joints of the same type across robots; the core aggregates the whole
body, and its gradients agree between robots with similar bodies [X] (Sec.~\ref{sec:where}). Sharing should
therefore follow the joint in the joint-level maps and the embodiment in the embodiment-level map.

Sparse mixtures of experts (MoE) meet the first requirement. A router sends each input to a few experts from a shared
pool, so that two inputs update common parameters only through the experts they both
select~\citep{shazeer2017outrageously,fedus2022switch}, and expert layers have been used to reduce the interference
between tasks~\citep{chen2023modsquad,hendawy2024moore}. In these models, however, the routed unit is a token or a
task. A cross-embodiment policy leaves the second requirement open: which unit to route at each stage of the policy,
and which morphological relationships the routing should follow.

To this end, we introduce Morphology-aware Routing of Shared Experts (MORSE), a framework that assigns the updates of
dissimilar joints and embodiments to different parameters by routing every stage of the policy at the level of its
computation. Each expert layer combines an always-active shared expert, which retains the computation common to all
robots, with a pool of routed experts shared across robots. In the joint-level layers, which encode the observation
and predict the action of each joint, a router conditioned on the properties and the state of the joint selects a few
experts for every joint. In the embodiment-level layers, which process the aggregated body representation, a router
conditioned on the embodiment and its state selects experts for the whole robot, and the result informs the action of
every joint. A mutual-information objective aligned with these two levels encourages joints of the same type and
robots of the same embodiment to rely on the same experts. Apart from this routing objective, MORSE leaves the offline
RL algorithm unchanged, so it applies to standard offline RL algorithms and complements methods that change how the
policy is updated. In summary, the main contributions of this paper are threefold:
\begin{itemize}
    \item We analyze where the updates of a cross-embodiment policy interfere, and show that the interference follows
    the type of the joint in the joint-level maps and the embodiment in the core, and that grouping the update by robot
    leaves it unchanged, since all groups update the same parameters.
    \item We propose MORSE, which assigns interfering updates to different experts and routes every stage of the policy
    at the level of its computation, every joint in the joint encoder and decoder and every embodiment in the core,
    with a routing objective that aligns the experts with groups of similar joints and with embodiments.
    \item We provide a thorough empirical study on 16 legged robots, six datasets, and three offline RL algorithms,
    showing that MORSE reduces the interference [X] and outperforms robot-level grouping and conflict-resolution
    methods, together with controlled comparisons of routing granularity, grouping, capacity, and robot reweighting.
\end{itemize}

% =====================================================================================================================
% PART B: NEW SECTION, inserted after Problem Formulation (Sec. 3) and before Method (Sec. 4)
% =====================================================================================================================

\section{Where Do the Updates Interfere?}
\label{sec:where}

Before presenting MORSE, we examine how joints and robots interact when they update the shared maps of
Eq.~\eqref{eq:policy}. We train the morphology-agnostic policy of~\citet{bohlinger2024onepolicy} with IQL on the 70\%
Suboptimal Forward dataset [X: setup] and, at checkpoints throughout training, compute the gradient of the actor loss
of Eq.~\eqref{eq:actor_loss} for individual units. For the joint-level maps $F$ and $D$, a unit is a joint of a robot,
and its gradient is the one that the map receives through the computation of that joint; for the embodiment-level map
$C$, a unit is a robot. As in~\citep{yu2020pcgrad,abe2026crossembodimentoffline}, two units conflict when the cosine
similarity of their gradients is negative. We compare joints by the joint groups of Sec.~\ref{sec:routing}, which
cluster the joints of all robots by their kinematic and dynamic properties, and robots by the morphology groups of EG.
Figure~\ref{fig:where} summarizes the results.

\begin{figure*}[t]
\centering
\fbox{\rule{0pt}{1.5in}\rule{0.95\textwidth}{0pt}}
\caption{Where the updates of a shared cross-embodiment policy interfere (IQL, 70\% Suboptimal Forward).
(a) Joint encoder and decoder: gradient cosine between joints of the same and of different joint groups, across robots
and within one robot. (b) Core: gradient cosine between robots, ordered by the groups of EG. (c) Fraction of
conflicting pairs within and between EG groups, under the joint update and during training with EG. [X]}
\label{fig:where}
\end{figure*}

\paragraph{Joint-level interference follows the type of the joint.}
In the joint encoder and the decoder, joints of the same group agree across robots: their gradients have a mean
cosine of [X], compared with [X] for joints of different groups, of which [X]\% conflict. The disagreement does not
stop at the boundary of a robot. Within one robot, joints of different groups conflict in [X]\% of pairs, as often as
joints of different robots [X]. A joint thus has more in common with the joints of the same type in other robots than
with the other joints of its own body, a relationship that no decision made per robot can express.

\paragraph{Embodiment-level interference follows the body.}
The core processes the aggregated representation of the whole body, and here the structure is that of robots: robots
with similar bodies have aligned gradients, with a mean cosine of [X] within EG groups compared with [X] between them.
At this level, the embodiment is the natural unit of sharing.

\paragraph{Grouping the update does not separate the updates.}
EG updates its groups in turn, and yet conflicting pairs are as frequent within its groups as between them ([X]\% and
[X]\%), and the fraction of conflicting pairs during training with EG stays close to that of the joint update ([X]\%
and [X]\%). This is expected: a grouped update changes when a robot writes to the parameters, not which parameters it
writes to. All groups update the same weights, and since they share the state of the optimizer, the step taken for one
group is dominated by the momentum of the others; with the first-moment rate $\beta_1=0.9$ of Adam, the gradient of the
current group enters the step with weight $1-\beta_1=0.1$.

\paragraph{Implications.}
Two requirements follow. Interfering updates must reach different parameters, which calls for parameters that are not
shared by all units, and the unit that decides the sharing must match the computation: the joint in the joint-level
maps and the embodiment in the core. MORSE meets both with sparse expert layers routed at the level of each map
(Sec.~\ref{sec:method}).

% =====================================================================================================================
% ANALYSES THAT FILL THE [X] (not part of the paper)
% =====================================================================================================================
% A1  Joint level, same vs different group (revised k=8 table): mean cosine and % conflicting, cross-robot and
%     within-robot pairs, per checkpoint, 2 seeds. Recomputable from the saved dense-URMA gradient files
%     (experiments/analysis/results/conflict/data, 2k-100k, seeds 0/1) with conflict_grouping.py and the new table;
%     the within-robot, different-group split needs a small addition to the script. No training needed.
% A2  Core, within vs between EG groups: mean cosine. Same files, no training needed.
% A3  Conflicting pairs during EG training: port the conflict hook (morse/conflict.py) to uni_iql_eg and run it on EG
%     checkpoints (2k-50k, >=2 seeds); compare with the dense run at the same checkpoints.
% A4  Contribution 3, "MORSE reduces the interference": the same measure on MORSE checkpoints, on the shared expert
%     and on the routed experts (only pairs that select a common expert interact), against the dense policy and EG.
%     Add a control with random routing of the same sparsity, so that the reduction is not attributed to sparsity
%     alone, and relate the per-robot change in interference to the per-robot change in return.
